\chapter{Native Modern GPU APIs}
\label{ch:native-modern-gpu}
\label{ch:vulkan-backend}

Four public identities belong in this cluster, but they do not share one implementation.
They share an explicit-API design vocabulary: pipelines and resource transitions are visible,
shader delivery is a first-class build concern, and a successful device creation says little
about which CNA features the path actually proves.

\begin{longtable}{p{0.14\textwidth}p{0.21\textwidth}p{0.17\textwidth}p{0.13\textwidth}p{0.20\textwidth}}
\toprule
Identity & Extended draw route & RT2D / MRT & Query & Shader boundary \\
\midrule
\endfirsthead
\toprule
Identity & Extended draw route & RT2D / MRT & Query & Shader boundary \\
\midrule
\endhead
\texttt{VULKAN} & native & yes / yes & yes & CNA GLSL compiled to SPIR-V \\
\texttt{SDL\_GPU} & native with bounded colored tail & yes / yes & no real query & API-specific precompiled bytecode \\
\texttt{WEBGPU} & native with bounded colored tail & yes / one attachment only & no real query & WGSL and native WebGPU binding model \\
\texttt{METAL} & native & yes / no & no & embedded MSL compiled at runtime \\
\bottomrule
\end{longtable}

The table states behavior, not merely capability bits.  \texttt{SDL\_GPU} inherits the base
null query while the default capability policy reports support.  \texttt{WEBGPU} likewise
reports MRT and query support, but rejects more than one target and has no real query object.
Those are contract defects tracked by the audit, not reasons to promote a no-throw smoke test
to feature evidence.

Vulkan is CNA's second-most-mature renderer by the project's own account, and the renderer
whose history best illustrates a recurring theme in this Part: a class of bug where an
interface method accepts every parameter it is supposed to, but the implementation behind it
quietly discards most of them.

\section{Recorded Vulkan campaign baseline}

The recorded Vulkan validation campaign reported 4{,}371 passes in a 4{,}373-case configured
population, with two hardware-dependent skips; its renderer-specific \texttt{ctest} population
passed 126 of 127.  These are campaign populations, not a universal current test count; the
counting model is in Chapter~\ref{ch:counting-discipline}.
The single remaining failure, \texttt{Vulkan\_DepthBias}, is a narrow one: only the most
extreme tested magnitude of depth bias produces no observable effect
(\S\ref{sec:vulkan-depthbias}), a symptom shared with an equivalent finding on
\texttt{D3D9}. Historically this renderer also carried five failing \cnaclass{BlendState}
tests and three segfaulting skinned-model content tests --- both fully resolved, with no
test exclusions needed to reach the otherwise clean campaign baseline.

\section{BlendState: from ``almost entirely fake'' to fully real}

The single most serious historical bug on this renderer, in the project's own words, is that
its blend-state implementation was ``almost entirely fake'': one hardcoded blend equation was
applied regardless of what \cnaclass{BlendState} actually requested, confirmed failing across
five independent pixel tests before the fix. The fix replaced this with a genuine
per-\cnaclass{Blend} / \cnaclass{BlendFunction} mapping, applied consistently across all nine
of this renderer's own pipeline-creation call sites --- meaning the fix had to be threaded
through every place a graphics pipeline object gets built, not just one central function.

\section{Stencil testing: the same ``discarded parameters'' bug shape}

\cnaclass{DepthStencilState} (Chapter~\ref{ch:state-objects}) had an almost identical
history. The renderer's own state-application entry point accepted fifteen parameters but
genuinely stored only two of them (the depth-enable and depth-write flags) --- every
stencil-related parameter passed through the call was simply discarded, and the pipeline's
own stencil-test-enable flag was never set at all, on any pipeline. Five separate,
independently-run checks (the enable flag itself, the stencil masks, front-face operations,
two-sided stencil mode, and reference-stencil propagation) each separately confirmed this
failure mode. The fix added real per-pipeline depth-compare operations, full front- and
back-face \texttt{VkStencilOpState} configuration, and dynamic-state reference/mask values ---
and, as a direct side effect of the same fix, connected \texttt{ReferenceStencil}'s
independent-override path (Chapter~\ref{ch:state-objects}) on Vulkan specifically, something
still broken on EasyGL and BGFX. A second, smaller bug found in the same investigation: depth
format selection tried a stencil-\emph{less} format before it tried any stencil-capable
alternative.

\section{OcclusionQuery: from never-wired to a genuinely discriminating query}

Before its fix, occlusion queries on this renderer were not merely inaccurate --- they were
never connected to anything at all. This renderer defers 3D and 2D draws into a pending-batch
snapshot, recorded into real Vulkan commands only once per frame, well after
\texttt{Begin()} / \texttt{End()} had already returned synchronously; queries therefore always
reported zero, regardless of what was actually visible. The fix required resolving three
genuinely non-obvious design questions: how to tag an individual pending draw with the query
that should observe it; whether a single query may legitimately span multiple draw calls
(resolved: yes, as long as they share one render pass, tracked via contiguous-run detection,
with an explicit, documented limitation that a query spanning a render-pass \emph{boundary}
is not correctly summed --- a real capability gap against what a real device could do, stated
plainly rather than silently assumed); and how per-frame query-pool resets need to be
sequenced before any render pass begins. Verified in both directions --- a visible quad
producing a positive pixel count, an occluded quad behind a nearer opaque object producing
none --- plus a genuine multi-draw-span case (two non-overlapping half-quads inside one
\texttt{Begin} / \texttt{End} span, summing correctly), discriminating for real in this
project's own software Vulkan implementation.

\subsection{A worked example: using an OcclusionQuery correctly on this renderer}

The real \cnaclass{OcclusionQuery} surface is small --- \texttt{Begin()}, \texttt{End()},
\texttt{getIsCompleteProperty()}, \texttt{getPixelCountProperty()} --- and the standard
visibility-culling pattern is exactly what the verification above tested directly:

\begin{lstlisting}
OcclusionQuery query(getGraphicsDeviceProperty());

query.Begin();
DrawBoundingProxyGeometry(occluderCandidate); // a cheap stand-in, not the real detailed mesh
query.End();

// One or more frames later, once rendering has actually reached the GPU:
if (query.getIsCompleteProperty() && query.getPixelCountProperty() > 0) {
    DrawFullDetailMesh(occluderCandidate); // visible -- worth the real draw cost
}
\end{lstlisting}

The one genuine gotcha this renderer's own verification surfaced, worth restating as a rule
rather than only as a finding: a single query's \texttt{Begin()} / \texttt{End()} span may
legitimately cover \emph{multiple} draw calls (useful for querying a whole cluster of small
objects with one query instead of one per object), but only as long as every draw in that span
stays within one render pass. A query spanning a render-pass boundary --- switching render
targets, or anything else that closes and reopens a Vulkan render pass mid-span --- is not
correctly summed on this renderer today. A conservative rule that avoids the gap entirely:
never let a render-target switch happen between an \cnaclass{OcclusionQuery}'s \texttt{Begin()}
and \texttt{End()}.

\section{Remaining limitation}
\label{sec:vulkan-depthbias}

Exactly one genuine, unresolved gap remains on this renderer at the pinned revision: an
isolated \cnaclass{RasterizerState.DepthBias} extreme-magnitude sub-case, where real hardware
bias produced no observable effect at any tested magnitude in this project's own sandboxed
environment --- suspected to be an environment or driver limitation shared with an equivalent
finding made independently on \texttt{D3D9} (Chapter~\ref{ch:direct3d-backends}), rather than
a Vulkan-specific defect.

\section{Two distinct render-target black-render bugs, both now closed}
\label{sec:vulkan-rendertarget-black}

The repository's still-stale \texttt{docs/rendertarget-support.md} describes two confirmed
black-render defects as open. The live plan, source, and tests close both. Preserving the
mechanisms still matters because the first is part of the current clear architecture and the
second is a useful lesson in re-verifying an old diagnosis.

Task~875 was not specific to cube targets. \texttt{VulkanRenderer::Clear()} once changed
only global clear scalars, while render passes were collected solely from real draw calls. A
clear-only target therefore never entered command recording:

\begin{lstlisting}
device.SetRenderTarget(&probe, CubeMapFace::PositiveZ);
device.Clear(Color::CornflowerBlue);
device.SetRenderTarget(nullptr); // no intervening draw: the former Task 875 trigger
\end{lstlisting}

The fix adds the active target to a deduplicated \texttt{clearedRTs\_} list from both
\texttt{Clear()} and \texttt{ClearColorAndDepth()}; command recording seeds its used-target
list from there before scanning queued SpriteBatch and 3D draws. The dedicated clear-only
\cnaclass{RenderTarget2D} regression fills two targets without a draw, forces the necessary
frame boundaries, then samples both. Disabling both insertions
reproduced black; restoring them returned the exact fills. One important limit remains:
\texttt{clearedRTs\_} guarantees that a pass exists, but does not make the pass's attachment
load operations reflect each \cnaclass{ClearOptions} subset
(\S\ref{sec:backend-clear-matrix}).

Task~876 was the separate all-six-faces path: real SpriteBatch content was rendered into a
\cnaclass{RenderTargetCube}, unbound, then sampled through \cnaclass{EnvironmentMapEffect}.
The same original regression now produces the exact expected blue centre across five repeated
runs plus its registered test invocation. No weakened replacement test was substituted.
The project did not retroactively bisect which intervening change fixed it; the substantial
Task~907 cube-view/layout rework is the leading inference, but its most obvious transition fix
is mip-count-gated while this regression uses one level. The honest current conclusion is
therefore ``verified closed, precise historical fixing commit unknown,'' not ``still open.''

\section{Clear values are real; \texttt{ClearOptions} selectivity is not}

Task~871's cross-frame tests establish that colour, depth, and stencil values can all reach a
real Vulkan attachment. They do not make \cnaclass{ClearOptions} an issue-ordered command.
The backbuffer render pass always uses \texttt{LOAD\_OP\_CLEAR} for every existing attachment,
and a discard-content render target does the same. All calls made before queued draws merely
replace global clear scalars; the final values apply once at pass start, before every draw,
regardless of call order.

Preserve-content targets expose a different failure: their render pass fixes colour to
\texttt{LOAD} and depth/stencil to \texttt{DONT\_CARE}. Adding the target to
\texttt{clearedRTs\_} guarantees a pass will exist, but an explicit clear does not select a
different render-pass variant. \texttt{ClearDepth()} is narrower still: it changes only the
global depth scalar and does not add an otherwise-unused current target to that list. These
limits are separate from Task~875's now-fixed ``no pass at all'' defect and are tabulated
against the other audited renderer families in \S\ref{sec:backend-clear-matrix}.

The constructor comparison creates two additional usage boundaries. Only exact
\texttt{Preserve\allowbreak{}Contents} reaches the renderer as true, so
\texttt{Platform\allowbreak{}Contents} selects the discard pass even though shared code issues
no matching black clear. The MSAA pass has no load variant at all, so Preserve+MSAA also
discards. Cube targets cannot convey usage through the common factory, and MRT uses a separate
always-clear pass; both discard for every enum. The non-MSAA two-frame usage test proves only
the color Discard/Preserve split, not these depth, MSAA, cube, MRT, or Platform cases. See
\S\ref{sec:backend-rendertarget-usage-matrix}.

\section{A confirmed non-gap: device-feature gating}

Not every investigation into this renderer found a bug. A dedicated check of whether optional
device features (\texttt{fillModeNonSolid}, needed for wireframe rendering, and
\texttt{samplerAnisotropy}) are properly negotiated before being requested found the gating
already correctly implemented --- both are queried via \texttt{vkGetPhysicalDeviceFeatures}
before being enabled, with a graceful fallback if the underlying device lacks either one.
This is included here deliberately: a book that only ever reports bugs found would give a
skewed picture of a renderer whose overall health is, in fact, good.

\section{Notable feature-matrix specifics}

A cube-map render target's MSAA support is real but indirect: rather than a genuinely
multisampled cube attachment, it is implemented as an ordinary (non-cube) multisampled 2D
texture array paired with a separately cube-flagged resolve texture. Multiple render targets
are capped at four by shared, cross-renderer code mirroring FNA's own
\path{MAX_RENDERTARGET_BINDINGS} constant --- notably \emph{not} because Vulkan's own
device capability is the limiting factor, since a typical device supports eight. Uploading
\texttt{Texture2D} data at a mip level above zero is now real: the image and view allocate
the public level count, a level-specific barrier and staged copy update the requested
subresource, and the dedicated minification test samples its authored color. The former
silent no-op recorded earlier in this chapter's source history was fixed by Task~925; see
Chapter~\ref{ch:backend-architecture}, \S\ref{sec:texture2d-update-matrix}, for the current
cross-renderer matrix and proof boundary. Forwarding a real, non-\texttt{Color}
\cnaclass{SurfaceFormat} shares the identical blocked/architectural
limitation already covered for EasyGL in Chapter~\ref{ch:easygl-backend} --- both renderers
are blocked by the same upstream \texttt{Texture::ValidateFormat} contract, not by anything
Vulkan-specific.

\section{Three tracked gaps, checked against live source rather than trusted from the doc}

This project's own \texttt{docs/rendertarget-support.md} --- the same document
\S\ref{sec:vulkan-rendertarget-black} already drew Tasks~875/876 from --- lists three further
tracked tasks for this renderer as still open: real GPU mip generation for a render target
(Task~878), real render-target MSAA (Task~879), and any actual GPU effect from a custom
sub-region \cnaclass{Viewport} at all (Task~880, ``zero wiring on any renderer''). Reading
\texttt{VulkanRenderer.cpp} directly, rather than trusting that document's own status
column, finds all three already resolved --- each superseded by a later task number
(\texttt{Task~903} / \texttt{907}) whose own comments the doc was never updated to reflect, one
more real instance of the recurring ``a status that looked settled turned out to be stale''
pattern this book has run into more than once. \texttt{VulkanRenderTargetRenderer::%
MaybeGenerateMips()} runs a real, per-level \texttt{vkCmdBlitImage} cascade against a render
target's just-rendered (and, if applicable, just-MSAA-resolved) level-0 content immediately
after its render pass ends --- the direct Vulkan equivalent of EasyGL's own
\texttt{glGenerateMipmap}-on-unbind. Render-target MSAA is real for both
\cnaclass{RenderTarget2D} and \cnaclass{RenderTargetCube} (the cube case exactly as described
above). And \texttt{SetViewport()} genuinely reaches \texttt{vkCmdSetViewport} at
command-buffer-record time, mirroring \texttt{SetScissorRect}'s own identical storage-then-apply
shape --- with one honestly-scoped limitation worth stating precisely rather than rounding up to
``fixed'': a custom sub-region \cnaclass{Viewport} is only honored for the backbuffer pass. A
render-target pass stays hardcoded to that target's own full size, because this renderer's
deferred, potentially-multiple-render-targets-per-frame recording model has no way to recover
\emph{which} \cnaclass{Viewport} was actually active at the moment each render target's own draws
were originally issued, by the time a frame's commands actually get recorded.

The same scope limit applies to scissor, not only viewport. Each render-target pass hardcodes
both a full-target \texttt{VkViewport} and \texttt{VkRect2D}; the stored custom values are read
only for the backbuffer pass, with \texttt{ScissorTestEnable=false} or a zero rectangle falling
back to the full swapchain. One dynamic viewport/scissor pair is installed before both
\texttt{drawSpritesFor()} and \texttt{draw3DFor()}, so the last live state when the command
buffer is recorded governs every queued sprite and 3D draw in that backbuffer pass. It is not
snapshotted per draw at the time the game issued it. Dedicated pixel tests prove the custom
backbuffer viewport and enabled/disabled scissor cases; the shared render-target reset test
proves the full-size reset visible after switching targets, not a custom sub-region inside a
Vulkan target.

\section{How a swapchain gets built: format, color space, and present mode}
\label{sec:vulkan-swapchain}

\texttt{VulkanRenderer::CreateSwapchain()} is where two real, easy-to-get-wrong
decisions get made once, at swapchain-creation time, rather than left to chance: which pixel
format the swapchain presents through, and which \texttt{VkPresentModeKHR} actually governs
frame pacing. Both are worth reading directly rather than assumed, since Vulkan's own API
surface makes it easy to pick a format that looks equivalent but silently changes what a game
sees on screen.

\subsection{UNORM, not SRGB: a deliberate color-space decision}

A physical device typically reports several supported \texttt{(VkFormat, VkColorSpaceKHR)}
swapchain pairs, and it is common in Vulkan tutorials to prefer an SRGB-variant format (e.g.\
\texttt{VK\_FORMAT\_B8G8R8A8\_SRGB}) for its automatic linear-to-sRGB gamma encode on every
presented pixel. This renderer deliberately does the opposite, and says why directly in its own
source comment: real XNA's default backbuffer format, \cnans{SurfaceFormat}\texttt{::Color}
(Appendix~\ref{ch:api-quick-reference}), is linear --- a separate,
CNAEXT\allowbreak{} \texttt{ColorSrgbEXT} value exists specifically for the gamma-encoded
variant (Appendix~E) --- so an SRGB swapchain format would apply an automatic gamma encode XNA
itself never performed, silently darkening every frame relative to what a real XNA game
produced. The fix is a preference search rather than a hardcoded pick, since not every device
lists the exact pair in the same order:

\begin{lstlisting}
VkSurfaceFormatKHR fmt = fmts[0];  // fallback: whatever the device lists first
for (auto& f : fmts)
    if (f.format == VK_FORMAT_B8G8R8A8_UNORM &&
        f.colorSpace == VK_COLOR_SPACE_SRGB_NONLINEAR_KHR)
    { fmt = f; break; }
\end{lstlisting}

The two fields answer two different questions, which is exactly the subtlety worth stating
plainly: \texttt{colorSpace} here is \texttt{VK\_COLOR\_SPACE\_SRGB\_NONLINEAR\_KHR} regardless
--- that value describes the \emph{display}'s own expected color space (the near-universal
default virtually every consumer monitor and every Vulkan implementation supports), not
whether the swapchain itself gamma-encodes on write. \texttt{format} is the field that
controls the encode: \texttt{\_UNORM} writes the raw bytes a game submits with no
transformation, while a sibling \texttt{\_SRGB} format of the identical byte layout would
gamma-encode automatically. Picking \texttt{UNORM} with the standard nonlinear color space is
therefore not a contradiction --- it is precisely ``present to a normal monitor, but do not
add an encode step XNA's own \texttt{Color} format never had.'' A porting engineer who
free-associates ``SRGB'' with ``correct color space, always prefer it'' and swaps this one
line risks introducing a real, visible darkening regression across every scene in the game,
not just an edge case.

\begin{sourcenote}
That deliberate UNORM choice is fixed renderer policy, not a response to the caller's current
\texttt{BackBufferFormat}. The Vulkan factory discards both presentation format ordinals:
it prefers this same B8G8R8A8 UNORM pair even when public
\cnaclass{PresentationParameters} reports something else, and always creates a
device-selected depth image (D24S8, then D32S8, then D32) even for
\texttt{DepthFormat::None}. Render-target depth formats use a newer, genuinely per-target
selection path and must not be confused with this default-backbuffer behavior. SDL can still
change the window to fullscreen; a later out-of-date/resize rebuild follows the surface size
without consuming \texttt{isFullScreen} itself.
\end{sourcenote}

\subsection{SwapInterval to VkPresentModeKHR: the same three-way switch Chapter~\ref{ch:game-loop} introduced}

\cnaclass{PresentationMode}'s \texttt{PresentInterval} (Immediate / Default-a.k.a.-One / Two,
Chapter~\ref{ch:game-loop}) is not merely stored on this renderer --- it selects among the
device's actually-supported present modes with an explicit fallback chain for each case,
verified directly against \texttt{CreateSwapchain()}:

\begin{lstlisting}
VkPresentModeKHR mode = VK_PRESENT_MODE_FIFO_KHR; // safe default: guaranteed available
if (swapInterval_ == 0) {         // Immediate: tearing allowed, lowest latency
    // prefer true IMMEDIATE; MAILBOX (triple-buffered, no tearing) is the fallback
    // if IMMEDIATE isn't supported; FIFO if neither is
} else if (swapInterval_ == 2) {  // Two: allow a slower-than-vsync present without stalling
    // prefer FIFO_RELAXED; falls through to FIFO if unsupported
}
// swapInterval_ == 1 (Default/One): FIFO unconditionally -- always available, real VSync
\end{lstlisting}

Two things about this chain are worth internalizing rather than skimming past. First,
\texttt{FIFO} is the only present mode the Vulkan specification \emph{guarantees} every
conformant implementation supports --- choosing it as the unconditional starting value, and as
every fallback's own final fallback, is what makes swapchain creation itself never fail for
lack of an exotic present mode, on any device. Second, \texttt{PresentInterval::Two}'s mapping
to \texttt{FIFO\_RELAXED} is a deliberately different request from \texttt{Immediate}'s: FIFO
Relaxed still waits for vsync when the application is keeping up, and only presents
immediately (accepting a single tear) on the specific frame where the application would
otherwise have missed vsync and stalled --- the Vulkan-native shape of ``prefer vsync, but
don't let a single slow frame cascade into a visible stutter,'' not a request for tearing as a
matter of course the way \texttt{Immediate} is. A game that sets \texttt{PresentInterval::Two}
expecting literally half the refresh rate (the historical, D3D9-era ``present interval two''
reading some porting engineers carry over) will not get that on this renderer --- it gets FIFO
Relaxed's own, different throughput/latency tradeoff instead, worth checking explicitly against
a real target frame-rate requirement before assuming interval semantics carried over unchanged
from a Direct3D background.

\begin{sourcenote}
That switch runs only in \texttt{CreateSwapchain()}, from the interval captured by the
renderer constructor. \texttt{VulkanGraphics\allowbreak Renderer} does not override
\texttt{IGraphicsRenderer::SetSwap\allowbreak Interval()}, so a later
\cnaclass{GraphicsDevice}\texttt{::Reset()} stores the new public value but cannot update
\texttt{swapInterval\_} or recreate the swapchain for it. This is especially visible through
\cnaclass{Game}: its device is eagerly built with default \texttt{PresentInterval::Default}
before the derived constructor can set
\texttt{GraphicsDevice\allowbreak Manager.SynchronizeWith\allowbreak VerticalRetrace=false}.
The subsequent
\texttt{ApplyChanges()} therefore still leaves Vulkan in FIFO mode. Only constructing a
\cnaclass{GraphicsDevice} directly with the desired parameters (or using CNA's separate
test-only full-renderer recreation helper) can currently change that initial choice. No
dedicated test queries or times the real Vulkan present mode.
\end{sourcenote}

\section[Two frames in flight, and the deferred-present trick]{Two frames in flight, and
the deferred-present trick that makes readback pixel-safe}
\label{sec:vulkan-frames-in-flight}

Everything in \S\ref{sec:vulkan-swapchain} covers how the swapchain itself gets built; this
section covers what actually happens every frame once it exists. This renderer double-buffers
its own CPU-side work against the GPU --- \texttt{MaxFramesInFlight} is a real, fixed
\texttt{constexpr int} value of~\texttt{2}, read directly from
\texttt{VulkanRenderer.hpp} --- with one \texttt{VkFence}, one pair of
\texttt{VkSemaphore}s (image-available, render-finished), and one command buffer per
in-flight frame, indexed by a wrapping \texttt{currentFrame\_} counter.

\subsection{The ordinary present path}

\texttt{Present()} calls the renderer's real \texttt{SubmitFrame(bool deferSwap)} with
\texttt{deferSwap=false}. Reading it directly shows five sequential steps, each with a real
purpose rather than boilerplate:

\begin{lstlisting}
vkWaitForFences(device_, 1, &inFlightFences_[currentFrame_], VK_TRUE, UINT64_MAX);
// ... vkAcquireNextImageKHR, vkResetFences, RecordCommandBuffer ...
vkQueueSubmit(graphicsQueue_, 1, &si, inFlightFences_[currentFrame_]);
// deferSwap == false: present immediately
vkQueuePresentKHR(presentQueue_, &pi);
currentFrame_ = (currentFrame_ + 1) % MaxFramesInFlight;
\end{lstlisting}

The opening \texttt{vkWaitForFences} is the throttle that makes double-buffering safe at all:
before this frame's slot can be reused, the GPU must have finished the \emph{previous} frame
that used the same slot two frames ago, or the CPU would start overwriting a command buffer
the GPU might still be reading. \texttt{vkAcquireNextImageKHR} signals
\texttt{imageAvailableSemaphores\_[currentFrame\_]} when the swapchain image is actually ready
to be written to (not necessarily immediately --- the presentation engine may still be
displaying it); \texttt{vkQueueSubmit} makes the GPU work wait on that same semaphore before
touching the image, and signals \texttt{renderFinishedSemaphores\_[currentFrame\_]} when done;
and \texttt{vkQueuePresentKHR} waits on that render-finished semaphore before handing the image
to the presentation engine. Three separate synchronization primitives, three separate
questions each one alone answers: is the CPU allowed to reuse this frame slot yet (fence), is
the GPU allowed to start rendering into this image yet (image-available semaphore), and is the
presentation engine allowed to display it yet (render-finished semaphore).

\subsection{The deferred path: why \texttt{GetBackBufferData} cannot simply call \texttt{Present}}

\cnaclass{GraphicsDevice}\texttt{::GetBackBufferData} (Chapter~\ref{ch:graphicsdevice}) needs a
real answer to a question the ordinary present path never has to ask: are the exact pixels the
GPU just rendered still readable, or has the presentation engine already taken ownership of
the image and possibly begun compositing it to the screen? Reading
\texttt{ReadBackbuffer()}'s own code comment states the race directly: \emph{``SubmitFrame(true)
renders and waits for the GPU but HOLDS the present, so we read the staging buffer below
BEFORE the image is handed to the presentation engine --- no race possible.''} The mechanism is
the same \texttt{SubmitFrame()} used for ordinary presentation, called with
\texttt{deferSwap=true} instead:

\begin{lstlisting}
if (deferSwap) {
    // Wait for render + readback copy to complete, but hold the image. The caller
    // (ReadBackbuffer) reads the staging buffer before the image is presented, so
    // presentation-engine timing can never corrupt the captured pixels.
    vkWaitForFences(device_, 1, &inFlightFences_[currentFrame_], VK_TRUE, UINT64_MAX);
    deferredPresentImageIndex_ = imageIndex;
    hasDeferredPresent_        = true;
    return true;
}
\end{lstlisting}

\texttt{RecordCommandBuffer} has already copied the whole swapchain image into a real staging
buffer (\texttt{readbackStagingBuf\_}) as part of the same command buffer the ordinary render
work goes into --- so this second \texttt{vkWaitForFences} call (on top of the one at the very
start of \texttt{SubmitFrame}) blocks until that copy is guaranteed complete on the GPU side,
before \texttt{ReadBackbuffer} maps \texttt{readbackStagingMem\_} and reads it directly.
\emph{Only after} the CPU has finished reading does \texttt{ReadBackbuffer} call
\texttt{FinishDeferredPresent()}, which performs the \texttt{vkQueuePresentKHR} call that
\texttt{SubmitFrame(true)} deliberately skipped --- late, but still exactly once, so no frame
is silently dropped from the screen because it was read back.

\subsection{Two further real findings from the same read path}

Reading \texttt{ReadBackbuffer()}'s own surrounding logic in full surfaces two further details
worth knowing, neither one hypothetical:

\begin{description}
  \item[A real read cache, not a re-render-every-call assumption.] A naive readback
        implementation might call \texttt{SubmitFrame(true)} unconditionally on every
        \texttt{GetBackBufferData} call --- but the real code guards it behind
        \texttt{hasNewWork \textbar\textbar{} !readbackStagingValid\_}, where
        \texttt{hasNewWork} checks whether any pending 3D or batched 2D draws actually exist
        since the last present. The code's own comment states why this matters concretely:
        re-presenting an already-rendered, empty-queue frame would re-render whatever
        \texttt{Clear} color is currently set and silently destroy the content every read
        after the first one would have returned, for any caller reading the same frame's
        backbuffer more than once (a golden-image test that reads several sub-regions of one
        frame, for instance).
  \item[\texttt{vkAcquireNextImageKHR} can genuinely fail on the very first frame.] A real,
        environment-specific comment in the same function notes that swapchain-out-of-date is
        \emph{``common on first frame under Wayland/RADV''} --- handled by zeroing the output
        buffer and returning, rather than throwing, so a caller can detect a blank frame and
        retry instead of crashing outright on a platform combination this specific to trigger.
\end{description}

A genuinely small but real detail rounds out the picture: swapchain images are frequently
\texttt{VK\_FORMAT\_B8G8R8A8\_UNORM} (per the color-space decision in
\S\ref{sec:vulkan-swapchain}), which stores channels in the opposite byte order from the RGBA
layout every other renderer and \cnaclass{Color} itself use --- \texttt{ReadBackbuffer} checks
\texttt{swapchainFormat\_} directly and performs a real per-pixel \texttt{B} / \texttt{R} channel
swap only when the format is actually BGRA, rather than assuming one fixed byte order across
every device this renderer might run on.

\section{A real, hittable ceiling: 512 live descriptor sets, shared across two resource types}
\label{sec:vulkan-descriptor-ceiling}

Every \cnaclass{Texture2D} and every \cnaclass{RenderTarget2D} constructed on this renderer
allocates one \texttt{VkDescriptorSet} from a single shared pool, sized at construction time
and never grown:

\begin{lstlisting}
static constexpr uint32_t MaxDescriptorSets = 512;
// ...
VkDescriptorPoolSize ps{ VK_DESCRIPTOR_TYPE_COMBINED_IMAGE_SAMPLER, MaxDescriptorSets };
VkDescriptorPoolCreateInfo ci{};
ci.flags = VK_DESCRIPTOR_POOL_CREATE_FREE_DESCRIPTOR_SET_BIT;
ci.maxSets = MaxDescriptorSets;
ci.poolSizeCount = 1;
ci.pPoolSizes = &ps;
vkCreateDescriptorPool(device_, &ci, nullptr, &descriptorPool_);
\end{lstlisting}

\cnaclass{VulkanTextureRenderer}'s own constructor and \cnaclass{VulkanRenderTargetRenderer}'s
own constructor both allocate from this identical \texttt{descriptorPool\_} --- confirmed by
reading both call sites directly, not just the pool's own creation code --- so the real budget
a game actually has is 512 \emph{simultaneously live} \cnaclass{Texture2D} plus
\cnaclass{RenderTarget2D} instances combined on this renderer, not 512 of each independently.
The \texttt{VK\_DESCRIPTOR\_POOL\_CREATE\_FREE\_DESCRIPTOR\_SET\_BIT} flag means this is a real
budget of \emph{concurrently undisposed} resources, not a lifetime total: each renderer's own
\texttt{ReleaseVulkanResources()} calls \texttt{vkFreeDescriptorSets} on \texttt{Dispose()},
returning its slot to the pool for reuse. A game that creates and disposes textures across many
frames --- streaming a level's assets in and out, for instance --- never accumulates toward the
ceiling as long as disposal genuinely keeps pace with creation.

The failure mode itself is a clean, loud exception rather than a silent cap or a corrupted
allocation, confirmed by reading the allocation call site directly:

\begin{lstlisting}
if (vkAllocateDescriptorSets(dev, &dsInfo, &descriptorSet_) != VK_SUCCESS)
    throw std::runtime_error("vkAllocateDescriptorSets (texture) failed");
\end{lstlisting}

A 513th simultaneously-live \cnaclass{Texture2D} or \cnaclass{RenderTarget2D} throws at
construction time, on this renderer specifically --- not a per-frame budget, and not shared with
any other renderer's own resource limits (EasyGL and BGFX have no equivalent fixed-pool
ceiling at all; each stock effect type on this same Vulkan renderer additionally has its own,
separately-sized dedicated pool --- \texttt{descriptorPoolEnvMap\_},
\texttt{descriptorPoolSkinned\_}, \texttt{descriptorPoolPbr\_}, and others --- so this specific
512-set ceiling applies only to plain \cnaclass{Texture2D} / \cnaclass{RenderTarget2D}
allocation, not to how many objects can simultaneously use a given stock effect). A game
targeting Vulkan specifically with a very large, non-atlased texture library, or one that
creates many short-lived render targets without disposing the previous batch first, is the
realistic way to reach this ceiling in practice --- worth knowing as a real, renderer-specific
constraint before assuming texture/render-target counts that work fine on EasyGL port
unchanged to Vulkan.

\section{A documented, not a hidden, coverage gap}

One honest gap in \emph{testing}, not rendering: at the pinned revision there is no dedicated
pixel test for \cnaclass{SpriteBatch}'s sort-mode/rotation/scale/crop/flip behavior, for
\cnaclass{SpriteFont} glyph placement, or for a multi-mesh \cnaclass{Model} hierarchy on this
renderer specifically --- confirmed by directly searching the test suite for matching test
names and finding none, despite basic \cnaclass{SpriteBatch}/2D-demo smoke tests passing.
This does not mean the underlying shared C\texttt{++} code is wrong (it is the same code
every other renderer runs), only that this specific renderer has not had these specific
features independently re-verified against its own rendering path.

\section{METAL: a narrow implementation with explicit refusals}

\texttt{METAL} is macOS-only, enables Objective-C++, and links Apple's system frameworks.
Its MSL lives as an embedded source string and is passed to
\texttt{newLibraryWithSource:} at runtime.  That gives it a simpler packaging story than the
SPIR-V pipelines, but also makes compiler and device creation part of runtime startup.

The implementation is intentionally transparent about its present boundary.  It has native
extended primitive submission and \cnaclass{RenderTarget2D}, but deliberately rejects MRT,
MSAA, custom effects, occlusion queries, instancing, multi-stream input, and backbuffer
readback.  These refusals cite missing macOS pixel evidence rather than pretending that
interface presence proves rendering correctness.

Its automated macOS path proves construction, clear/present, and capability reporting; it
does not yet establish representative 3D draw pixels.  That is a sharper statement than
calling the renderer either complete or untested.  The device path executes, while much of
the feature surface remains deliberately outside the claimed evidence boundary.

\subsection{Why explicit refusal is useful}

The renderer contract's default-true capability polarity makes accidental overclaiming easy.
Metal's policy demonstrates the safer integration pattern: enumerate the unsupported surface,
throw at the boundary, and widen it only when a macOS-specific observable proves the new path.
Callers get an early diagnostic instead of a later blank frame.

\section{Comparing the four evidence tiers}

\texttt{VULKAN} has the broadest renderer-specific pixel history in this group.
\texttt{SDL\_GPU} is exercised through its currently selected native driver, which is not
proof of every driver SDL can choose.  \texttt{WEBGPU} has real application and pixel
evidence but retains the MRT/query capability overclaims above.  \texttt{METAL} has a real
macOS lifecycle gate while its 3D pixel surface remains explicitly unproven.

For portability work, record both axes: which CNA identity ran and which native driver or
device path served it.  An abstraction selecting Vulkan on one machine and Metal on another
does not turn one successful run into proof of both.

% The WebGPU and SDL GPU deep dives remain separate source files but compile as sections here.
\section{Legacy WebGPU Chapter Material}
\label{ch:webgpu-backend}

\texttt{WEBGPU} is CNA's fifth graphics renderer, and its youngest: the project owner
explicitly lifted a former prohibition on this renderer and authorized its implementation.
It targets native \texttt{wgpu-native} (pinned to version 29.0.1.1), selected the same way
as every other renderer via \texttt{-DCNA\_GRAPHICS\_RENDERER=WEBGPU}. Its own tracking plan
spans fourteen phases and over 130 individually numbered tasks --- large enough that this
chapter covers its trajectory rather than every task.

\section{What ``experimental'' concretely means here}

Unlike the four renderers covered in the previous three chapters, \texttt{WEBGPU} is
deliberately excluded from the cross-renderer feature matrix used throughout
Chapters~\ref{ch:sdl-renderer}--\ref{ch:vulkan-backend}, on the matrix's own stated grounds
that its feature surface was not yet broad enough for a meaningful side-by-side comparison at
the time that matrix was written. This chapter instead draws on the renderer's own dedicated
documentation and milestone history.

\section{The native 2D baseline}

Early milestones established device/surface setup, clear/present, a working
\cnaclass{Texture2D} path, vertex/index buffer uploads, and a WGSL-based
\cnaclass{SpriteBatch} implementation. A 120-frame clear-and-present lifecycle test, and a
manually-reviewed \cnaclass{SpriteBatch} validation scene (covering cropping, tint/alpha,
rotation, flips, filtering, and all three address modes), both passed before the native 2D
baseline was declared closed. The most instructive milestone here, though, is an independent
one: a real, independently-developed XNA-style game (a \emph{Windows Phone Speedy Blupi}
port, from the \texttt{mobile-eggbert} sibling repository) was run against this renderer on a
real desktop, reached its main menu, rendered its \cnaclass{SpriteBatch} content
pixel-correctly, and played through a real animated cutscene --- with progress bar and
cross-fade --- triggered by a simulated button click, with no WebGPU validation errors of any
kind. This is a genuinely different, and stronger, kind of proof than any of this renderer's
own purpose-built test scenes: an unmodified, independently-authored game simply working.

\section{A bug a manual screenshot review missed, and a pixel assertion caught}

One specific finding is worth naming as a methodological lesson in its own right: the
SpriteBatch pipeline's blend factors did not actually match the non-premultiplied output the
shader itself produced, so a translucent sprite --- one with an alpha value strictly between
fully transparent and fully opaque --- rendered fully opaque instead. This was completely
invisible to the earlier manual screenshot review covering rotation, flips, and filtering; it
was only caught once a real pixel-value assertion replaced ``does this look right in a
screenshot.'' The fix matched Vulkan's own blend-factor pairing. This is the same lesson
Chapter~\ref{ch:spritebatch} drew from SDL\_Renderer's ignored-transform-matrix bug: a
rendering pipeline can look entirely correct in every way a human screenshot review checks,
and still be measurably wrong.

\section{3D depth testing and the gap it exposed}

The first 3D draw path (\texttt{DrawColoredPrimitives} / \texttt{DrawIndexedColoredPrimitives})
was verified with a genuine near/far depth-ordering test rather than trusting draw order
alone. Building it surfaced a real, previously-invisible gap: applying a
\cnaclass{DepthStencilState} had no effect on this renderer at all before this point --- the
apply-state entry point was entirely unimplemented, meaning
\texttt{GraphicsDevice.DepthStencilState} had silently done nothing on WebGPU up to that
point. From there, \cnaclass{PbrEffect} (a real glTF~2.0 metallic-roughness BRDF, ported line
for line from EasyGL's own GLSL implementation) and skinned variants of both
\cnaclass{SkinnedEffect} and \texttt{PbrEffect} (bone-palette skinning up to 72 bones, again
ported directly from EasyGL) were added and both pass their dedicated test suites in full.

\subsection{PbrEffect: a CNAEXT metallic-roughness material}

\cnaclass{PbrEffect} is worth a concrete look before moving on, since --- unlike every stock
effect in Chapter~\ref{ch:stock-effects} --- it has no real-XNA counterpart at all: real XNA
4.0 predates the glTF~2.0 metallic-roughness PBR model by years, so this is a \texttt{CNAEXT}
addition ported line for line from EasyGL's own GLSL implementation, not a port of anything
Microsoft ever shipped. Its property surface mixes ordinary \cnaclass{IEffectMatrices}/%
\cnaclass{IEffectLights} / \cnaclass{IEffectFog} plumbing (identical in shape to every stock
effect) with a genuinely new four-texture-map PBR material:

\begin{lstlisting}
PbrEffect effect(getGraphicsDeviceProperty());
effect.setWorldProperty(worldMatrix);
effect.setViewProperty(camera.View());
effect.setProjectionProperty(camera.Projection());

effect.setTextureProperty(&baseColorMap);               // glTF "baseColorTexture"
effect.setNormalMapProperty(&normalMap);
effect.setMetallicRoughnessMapProperty(&metalRoughMap); // one texture: G=rough, B=metal
effect.setOcclusionMapProperty(&aoMap);
effect.setEmissiveMapProperty(&emissiveMap);

// Both factors multiply the map's channel value (glTF's own "value * factor" convention),
// or stand alone with no map bound:
effect.setMetallicFactorProperty(1.0f);
effect.setRoughnessFactorProperty(1.0f);
effect.setEmissiveFactorProperty(Vector3(0.0f, 0.0f, 0.0f));

effect.setLightingEnabledProperty(true);
effect.EnableDefaultLighting();
\end{lstlisting}

The \texttt{MetallicRoughnessMap} single-texture, two-channel packing above is not an
implementation shortcut CNA introduced --- it is glTF~2.0's own real specified layout for this
material model (green channel roughness, blue channel metallic, red and alpha unused), and
\texttt{PbrEffect} follows it exactly so a metallic-roughness texture exported by any
standards-conformant glTF pipeline can be bound directly with no channel-repacking step. This
is also, concretely, the effect Chapter~\ref{ch:models-meshes}'s two independent skinning
systems both connect to on this renderer: the WebGPU port adds a skinned variant of
\texttt{PbrEffect} alongside skinned \cnaclass{SkinnedEffect}, so a glTF character rig
authored with a metallic-roughness material and driven by \cnaclass{AnimationPlayer} reaches
GPU skinning through this effect rather than through \cnaclass{SkinnedEffect} at all.

\section{Render targets, environment mapping, instancing, and mip generation}

\cnaclass{RenderTarget2D} support combines depth and stencil into a single format regardless
of what was requested (unlike Vulkan's newer per-render-target depth-format selection,
Chapter~\ref{ch:vulkan-backend}),
and needed a real type-safety fix when sampling a render target back as an ordinary texture.
MSAA support, once implemented, initially appeared to fail its own dedicated test --- but
investigation found the infrastructure was already correct, and the reported failure was a
defect in the \emph{test itself} (a missing \texttt{RasterizerState::CullNone}), not in the
MSAA implementation. \cnaclass{EnvironmentMapEffect} and genuine hardware instancing were
added together, with instancing implemented as a real per-instance vertex stream needing no
bind-group changes, unlike the effect families that do require one per draw.
\cnaclass{RenderTargetCube} support surfaced a cross-cutting sprite-geometry bug that has since
been fixed: the original queue path computed clip-space positions from the backbuffer's logical
dimensions even while a differently-sized off-screen target was bound. \texttt{REMED-GFX-019}
now derives dimensions from the current \cnaclass{RenderTarget2D} or cube face instead. Mip
generation for \cnaclass{Texture2D}
and \cnaclass{TextureCube} is implemented as a genuine filtered downsample (a series of
full-screen-triangle render passes, since \texttt{wgpu-native} v29 has no direct
blit-image-equivalent call). Notably, this renderer auto-regenerates the chain after every
level-zero write to a plain \cnaclass{Texture2D} or \cnaclass{TextureCube}. Real XNA/FNA and
the other audited plain-texture paths keep upper levels explicitly authored; SDL\_GPU is the
narrow exception for plain cubes, where a complete level-zero face upload also invokes its
native generator, while its plain 2D and 3D textures remain authored-level resources.
WebGPU's behavior prevents undefined upper levels after a level-zero upload, but it is a
deliberate timing divergence rather than an unqualified superiority: a later level-zero write
can overwrite upper levels that a game authored earlier. Render-target mip generation remains
the separate unbind/pass-finalization contract.

\begin{sourcenote}
The default surface is equally independent of public presentation-format fields, but its
choice is unusually visible. \texttt{ConfigureSurface()} prefers BGRA8-sRGB, then
RGBA8-sRGB, before either UNORM variant. Thus a public \texttt{Color} request can accompany
a native sRGB swapchain, the opposite default-color policy from Vulkan. Its fixed depth
texture uses \texttt{Depth24PlusStencil8}. The renderer recreates it for every non-minimized
surface, including when the public depth value is \texttt{None}. SDL performs any fullscreen window transition;
the next frame detects its physical size and reconfigures the surface, but neither requested
format enum participates.
\end{sourcenote}

\subsection{The corrected render-target-size trap}

The former sprite-geometry bug is worth walking through because it produced no error of any
kind --- only a subtly wrong result:

\begin{lstlisting}
// e.g. a 256x256 render target, while the game window itself is 1920x1080:
device.SetRenderTarget(&smallOffscreenTarget);
spriteBatch.Begin();
spriteBatch.Draw(sourceTexture, Vector2::Zero, Color::White); // "fill the target"
spriteBatch.End();
\end{lstlisting}

Before \texttt{REMED-GFX-019}, the same \texttt{Draw()} call intended for the
$256\times256$ target was converted as though the target were the $1920\times1080$ window,
silently changing its scale and placement. The current \texttt{QueueSprite()} instead selects
the bound 2D target's width/height, the bound cube face's size, or the backbuffer logical
viewport as three explicit cases before converting pixels to NDC.

The dedicated \texttt{WebGPU\_SpriteBatch\_RenderTarget} regression makes the distinction
discriminating rather than visual-only: a $96\times72$ backbuffer is paired with asymmetric
$48\times32$ and $64\times40$ targets; the same destination rectangle must occupy the same
target-relative pixels in both. It also covers target--backbuffer--target isolation, a
non-identity SpriteBatch transform, texture orientation, and a $48\times48$ cube face. Probe
pixels inside the old half-scale rectangle but outside the correct one must remain black. The
old practical prohibition on differently-sized targets is therefore no longer needed.

\subsection{RenderTargetCube and its former silent cast failure}

\cnaclass{RenderTargetCube} support (Task WEBGPU-114) landed after the plain
\cnaclass{RenderTarget2D} path above, and reused most of its design directly: one shared
6-array-layer \texttt{WGPUTexture} (the same layout \texttt{WebGPUTextureCubeRenderer} uses
for an ordinary sampled, readback-capable \cnaclass{TextureCube}), one shared
\texttt{size}$\times$\texttt{size}
depth/stencil texture reused across all six faces (only one face is ever bound at a time, the
same choice Vulkan's own \cnaclass{RenderTargetCube} renderer makes), and a
\texttt{WGPUTextureViewDimension\_2D} view per face for the render-pass color attachment
alongside one \texttt{WGPUTextureViewDimension\_Cube} view spanning all six layers for sampling
the whole thing back later.

That whole-cube sampling view is where a real, previously-invisible bug lived. Before this
task, \cnaclass{EnvironmentMapEffect}'s \texttt{envMap} field on this renderer was hardcoded to
the concrete type \texttt{const WebGPUTextureCubeRenderer*} --- so binding an ordinary
\cnaclass{TextureCube} as \texttt{EnvironmentMapEffect.EnvironmentMap} worked, but
binding a \cnaclass{RenderTargetCube} (a different concrete class implementing the same
\cnaclass{ITextureCubeRenderer} interface) did not: the hardcoded-type cast failed silently,
and the effect fell back to its own 1$\times$1 white cube default instead of throwing or
producing an obviously wrong image. The fix is a small shared interface, read directly from
the real header:

\begin{lstlisting}
// WEBGPU-114: the cube-map sibling of IWebGPUSamplable -- lets both
// WebGPUTextureCubeRenderer and WebGPURenderTargetCubeRenderer
// (render-into) resolve safely to their own whole-cube view wherever a
// GpuDrawParams::envMap needs to be sampled, via a dynamic_cast that
// degrades to nullptr for a null or incompatible-type input.
class IWebGPUCubeSamplable
{
public:
    virtual ~IWebGPUCubeSamplable() = default;
    [[nodiscard]] virtual WGPUTextureView CubeView() const = 0;
};
\end{lstlisting}

Both concrete classes now implement \texttt{IWebGPUCubeSamplable} alongside their own primary
interface --- \texttt{WebGPUTextureCubeRenderer} and \texttt{WebGPURenderTargetCubeRenderer}
alike --- and the draw-dispatch code resolves through it with a single \texttt{dynamic\_cast}
rather than the old hardcoded
concrete-type cast:

\begin{lstlisting}
[[nodiscard]] const IWebGPUCubeSamplable* ResolveCubeSamplable(const ITextureCubeRenderer* tex)
{
    return tex != nullptr ? dynamic_cast<const IWebGPUCubeSamplable*>(tex) : nullptr;
}
// ...
command.envMap = ResolveCubeSamplable(params.envMap);
\end{lstlisting}

\texttt{modules/renderers/webgpu/examples/\allowbreak{}webgpu\_rendertargetcube\_test.cpp}'s own Check C is built specifically to
prove this fix rather than merely exercise the feature: it clears all six faces of a real
\cnaclass{RenderTargetCube} to blue, binds it as \texttt{EnvironmentMapEffect.EnvironmentMap},
draws a full-screen reflective quad, and asserts the center pixel is genuinely blue --- a check
that could not have passed under the old hardcoded-type cast no matter how correctly the rest
of the render-target machinery worked, since the cast itself was the failure point.

The same test file's Check D is worth walking through on its own, because it targets a
different, easy-to-get-wrong question: does switching \emph{into} a cube face and back
correctly restore the backbuffer's own render pass, with nothing leaking either direction?

\begin{lstlisting}
device.SetRenderTarget(nullptr);              // targeting the backbuffer
device.Clear(Color::Blue);
device.SetRenderTarget(&cubeTarget, CubeMapFace::PositiveZ);
device.Clear(Color::Red);                     // targets the cube face, not the backbuffer
device.SetRenderTarget(nullptr);              // back to the backbuffer

// The backbuffer's own centre pixel must STILL be Blue -- the intervening
// cube-face Clear(Red) must not have leaked into it -- and the cube face
// itself must genuinely have received Red, not been silently skipped.
\end{lstlisting}

This mirrors the identical architecture check Task WEBGPU-53/54 already established for plain
\cnaclass{RenderTarget2D} in the section above, extended to confirm the same "eager flush on
target switch" design generalizes correctly to a third target kind rather than being
coincidentally correct for the two it was originally built against. Two scope cuts are
documented honestly rather than silently: mip-chain regeneration throws for
\cnaclass{RenderTargetCube} (\texttt{mipMap=true} is rejected, matching plain
\cnaclass{RenderTarget2D}'s own precedent on this renderer), and MSAA is not implemented ---
\texttt{MultiSampleCount} is accepted and ignored, always reporting \texttt{0} honestly rather
than silently pretending to honor a request it cannot fulfil.

\subsection{Why block-compressed texture upload remains open}

The ``no renderer anywhere in the project does real block-compressed texture upload'' gap named
below is worth a fuller account on this renderer specifically, because it was genuinely
investigated here rather than assumed unfixable. A direct diagnostic against this project's own
pinned \texttt{wgpu-native} v29.0.1.1, run on the real development GPU, confirmed the hardware
side is not the obstacle at all: \texttt{wgpuAdapterHasFeature(adapter,
TextureCompressionBC)} returned true, and requesting a device with that feature in
\texttt{requiredFeatures} succeeded too --- this machine's real adapter genuinely supports
DXT1/3/5 upload today. The actual blocker sits one layer up, and applies project-wide, not just
to \texttt{WEBGPU}: \texttt{Texture2D::FromStream()} always fully CPU-decompresses DXT-compressed
source bytes to plain RGBA8 \emph{before} any renderer ever sees them, and the shared
\texttt{CNA::Internal::Graphics::ImageData} struct every renderer's texture-creation path
consumes has no field at all for a surface format or a compressed-bytes flag --- it is
hardcoded to a plain \texttt{std::vector<uint8\_t>} of RGBA8 pixels. A WebGPU-only workaround
was considered and deliberately rejected: hand-crafting a direct \texttt{IGraphicsRenderer::%
CreateTexture()} call with raw BC1 bytes stuffed into that RGBA8-typed field would compile and
might even render correctly on this one renderer, but would be reachable by no real game (nothing
in the actual \cnaclass{Texture2D} loading path could ever produce it) and would misrepresent
\texttt{ImageData}'s own documented contract --- exactly the shape of shortcut this project's
own culture treats as worse than leaving the gap honestly open. Closing this for real needs a
cross-renderer \texttt{ImageData} design change (or a parallel compressed-texture entry point)
plus a matching \texttt{Texture2D.cpp} change to stop always-decompressing DXT source data ---
deliberately scoped out of any single renderer's own task list, tracked as a dedicated follow-up
instead.

\section{Three former limitations that are no longer true}

The WebGPU plan moved quickly after this chapter's first draft, so three earlier shorthand
limitations need a precise correction rather than being allowed to turn into folklore. First,
an arbitrary plain \cnaclass{Texture2D} now has a real renderer-level \texttt{GetData()} path.
That method on \texttt{WebGPUTextureRenderer} copies the requested mip level into a temporary
\texttt{MAP\_READ} buffer with WebGPU's required row alignment, waits through the asynchronous
map callback, and then extracts the requested $x,y,w,h$ rectangle. The dedicated
\texttt{WebGPU\_Texture2D\_GetData} test proves three independently useful cases: an exact
full-gradient round trip, a non-origin $2\times2$ sub-rectangle, and a distinct level-one mip
round trip. That is deliberately tested through \texttt{IGraphicsRenderer}, not merely through
the public \cnaclass{Texture2D}: the latter normally returns its shared CPU pixel shadow for a
plain texture, so an XNA-layer round trip alone could pass while this GPU copy path remained a
silent base-class no-op.

Second, culling, scissor, and viewport state are no longer merely stored values. Culling is
baked into the WebGPU 3D pipeline key and mapped to \texttt{WGPUCullMode}; a differential
pixel test proves that a deliberately chosen winding disappears under one XNA cull mode and
remains visible under the other. Scissor and viewport are genuine dynamic render-pass state via
\texttt{wgpuRenderPassEncoderSetScissorRect} and
\texttt{wgpuRenderPassEncoderSetViewport}. Wiring the latter exposed an otherwise-hidden
logical-versus-physical-size hazard: CNA's stored logical viewport can exceed the real
swapchain dimensions. The renderer therefore clamps it at application time, which is both a
WebGPU API requirement and what restored the pre-existing 2D and 3D regression tests.
Every queued SpriteBatch or 3D command captures its own viewport and scissor state at the public
draw call. Replay applies those snapshots dynamically through the two native setters on the
backbuffer, 2D-target, and cube-face paths; consecutive commands with identical state skip the
redundant native call. Thus several changes inside one deferred pass no longer collapse to the
last live value. The shared deferred pixel fixtures prove the per-command result, while the
WebGPU cardinality tests prove it did not require a pipeline variant, render-pass split, extra
submit, or one setter call for every same-state draw. A disabled or zero scissor expands to the
full current target instead of issuing an invalid zero-area rectangle.

Third, this is not a claim that every state object is complete. \cnaclass{FillMode}'s
\texttt{WireFrame} value is stored but polygon draws now refuse it because the pinned WebGPU API
offers no polygon-mode switch. The renderer stores the complete stencil configuration but has
deliberately not yet baked stencil operations
into its many pipeline families. The last deferral is evidence-based rather than cosmetic:
the cull-mode mapping itself required a pixel test to correct an initially plausible but wrong
winding interpretation, so copying a stencil front/back mapping without an equivalent
differential stencil test would risk making an invisible wrongness look ``implemented.''

\section{Clear ordering and usage are consistent across target types}

Each public clear enters the same ordered stream as all eleven deferred draw families, carrying
independent colour, depth, and stencil flags and values. Replay partitions one bind cycle into
native pass segments: leading/consecutive clears can fold into the next segment's load actions,
while a clear after a draw closes that segment so the following pass observes it at the exact
public position. A discard-content target may clear its first segment as its usage permits;
later explicit clears remain ordered rather than becoming one global final value. The shared
clear-options and ordered-clear fixtures cover backbuffer, 2D-target, cube-face, and MRT-rejection
boundaries without an intervening flush.

The shared FNA predicate maps only \texttt{DiscardContents} to discard; Preserve and Platform
both preserve. WebGPU receives that Boolean for 2D and cube targets, and the first native
segment selects clear versus load for colour, depth, and stencil together. Later segments load
prior contents unless an explicit ordered clear selects an aspect. A cube face owns its colour
view while all faces share the target's depth/stencil attachment, matching the reference
contract. Cube-usage and depth/stencil-usage regressions distinguish all three enums and the
six faces rather than inferring behavior from a combined clear.

The backbuffer's
\texttt{Presentation\allowbreak{}Parameters.Render\allowbreak{}Target\allowbreak{}Usage} remains
outside this off-screen policy, and plural binding still rejects MRT rather than creating more
than one attachment. The full comparison is \S\ref{sec:backend-clear-matrix}; the separate
usage matrix is \S\ref{sec:backend-rendertarget-usage-matrix}.

\section{What remains open}

Several gaps are shared with every other CNA renderer rather than unique to this one: no
renderer anywhere in the project does real block-compressed texture upload (the paragraph above
covers exactly why, and why it is a shared architectural gap rather than a per-renderer one),
and the cross-renderer design change that would make it reachable from a real game remains
larger than any one renderer task. Specific to this renderer: multiple simultaneous render
targets, actual stencil-operation mapping (not merely the already-stored state and dynamic
reference value), \texttt{WireFrame} rendering,
custom \cnaclass{SpriteBatch} effects, and custom WGSL effects authored outside the stock-effect
set. \cnaclass{BlendState} is no longer on that list for the one attachment WebGPU can bind:
all source/destination factors and operations, dynamic \texttt{BlendFactor}, slot-zero
\texttt{ColorWriteChannels}, and \texttt{MultiSampleMask} reach keyed SpriteBatch and 3D
pipelines with dedicated differential tests. Slots 1--3 remain inapplicable until MRT exists.
The two \cnaclass{RenderTargetCube} scope cuts also remain honest: mipmapped cube targets
throw at construction, and their per-target MSAA request is accepted but reports zero rather
than claiming an unimplemented multisampled array-texture resolve. Browser/Emscripten WebGPU is
entirely out of scope so far --- this renderer is native \texttt{wgpu-native} only.

\section{Verification methodology, and how it differs from D3D9's}

This renderer's own verification runs on a native automated CTest suite (gracefully skipping
where no Wayland/X11 display is available) plus a battery of dedicated pixel-assertion tests
per feature area, and the one independent real-game integration test described above. It is
explicitly \emph{not} covered by the oracle-versus-real-XNA methodology
Chapter~\ref{ch:direct3d-backends} describes for \texttt{D3D9} / \texttt{D3D11} / \texttt{D3D12}
--- that comparison against a genuine, running XNA~4.0 reference is reserved for the three
Direct3D renderers specifically, not applied here.

\section{Legacy SDL GPU Chapter Material}
\label{ch:sdlgpu-backend}

\texttt{SDL\_GPU} is a CNA graphics renderer, and one this book itself was slow to
give a proper account of --- Chapter~\ref{ch:roadmap} names that delay honestly as a gap in
this book, not in CNA's own code. This chapter closes it: a real, Linux/Vulkan-driver-verified
renderer implementing the full \cnaclass{IGraphicsRenderer} contract
(Chapter~\ref{ch:backend-architecture}), built on top of SDL3's own \texttt{SDL\_gpu} API
rather than a native vendor API of its own.

\section{Why this renderer costs nothing new}

\texttt{SDL\_gpu} dispatches internally to Vulkan, D3D12, or Metal depending on platform and
driver availability, and CNA's own \texttt{plan\_sdlgpu.md} states plainly why this renderer was
worth adding on top of the hardware-backed renderers the project already had: SDL3 is
already a vendored dependency of this project (\texttt{third\_party/SDL}), and
\texttt{SDL\_gpu.c} is already compiled into every prebuilt SDL3 package this repository
produces --- unlike Vulkan (the Vulkan SDK/loader) or WebGPU (\texttt{wgpu-native}, a separate
library), standing this renderer up required zero new third-party dependency, only new CNA-side
code against an API already sitting in the build. The main class is
\texttt{CNA::\allowbreak{}Internal::\allowbreak{}Renderers::\allowbreak{}SdlGpu::\allowbreak{}SdlGpuRenderer}
\index{CNA::Internal::Renderers::SdlGpu::SdlGpuRenderer@CNA::Internal::Renderers::\allowbreak{}SdlGpu::SdlGpuRenderer}.
Its source lives under
\texttt{modules/renderers/sdl-gpu/src/}. Build the target
\texttt{cna\_\allowbreak{}renderer\_\allowbreak{}sdl\_\allowbreak{}gpu}. Select this renderer with
\texttt{-DCNA\_\allowbreak{}GRAPHICS\_\allowbreak{}RENDERER=SDL\_\allowbreak{}GPU}.
The project's own naming table is explicit that this class name is deliberately distinct from
the pre-existing \cnaclass{SdlRenderer} in \texttt{SdlRenderer/} --- the older
\texttt{SDL\_Renderer} 2D API and the newer \texttt{SDL\_gpu} API are unrelated, and this project
was careful not to let the similar SDL naming collide.

\section{Precompiled bytecode only, in a fixed resource-binding order}

\texttt{SDL\_gpu} does not accept GLSL or HLSL source text at runtime the way
\cnaclass{ShaderEffect}'s own public contract promises a game --- \texttt{SDL\_CreateGPUShader}
takes only \texttt{SDL\_GPUShaderFormat}-tagged precompiled bytecode: SPIR-V for the Vulkan
driver, DXBC/DXIL for the D3D12 driver, MSL/\texttt{metallib} for the Metal driver. A further
constraint worth knowing before assuming any existing shader is reusable as-is: every stage's
samplers, storage textures, storage buffers, and uniform buffers must be declared in a fixed,
HLSL-style binding order (\texttt{t[n]} / \texttt{s[n]} in one address space, \texttt{u[n]} in
another, \texttt{b[n]} in a third) --- a different convention from the raw
\texttt{layout(set=, binding=)} numbering the existing \cnaclass{VulkanRenderer}'s own
GLSL shaders use. This means Vulkan's compiled SPIR-V is a useful \emph{algorithmic} reference
for this renderer's own shader authoring (the same lighting math, alpha-test logic, and
dual-texture blend) but not directly reusable output --- every shader's binding layout had to
be re-authored specifically for \texttt{SDL\_gpu}'s own convention, confirmed directly against
the real \texttt{.glsl} sources under this renderer's own
\texttt{modules/renderers/sdl-gpu/src/shaders/}
directory (23 files today, one pair per stock effect --- \texttt{colored3d},
\texttt{alpha\_test3d}, \texttt{dual\_texture3d}, \texttt{env\_map3d}, \texttt{lit\_textured3d},
\texttt{pbr3d} / \texttt{pbr\_skinned3d}, \texttt{skinned3d} / \texttt{skinned\_colored3d}, and more).

\subsection{A worked example: the Y-flip bug a real screenshot caught}

This renderer's own first real milestone --- a textured, tinted, rotated, multi-flip sprite
scene surviving a resize --- is a concrete instance of a theme this book returns to
repeatedly (Chapter~\ref{ch:testing-philosophy}): a ``did it throw'' test alone cannot catch a
wrong-but-plausible image. \texttt{SDL\_gpu}'s Vulkan driver flips clip-space Y internally, for
cross-renderer NDC consistency reasons of its own --- the vertex shader's first, naively-ported
NDC computation (\texttt{gl\_Position = vec4(ndc, 0, 1)}, the same shape every other renderer's
own sprite vertex shader already uses) compiled cleanly, ran without error, and rendered every
single sprite upside down, a texture's top row appearing at the bottom of the screen. Only a
real captured screenshot caught it; the fix negates Y explicitly
(\texttt{vec4(ndc.x, -ndc.y, 0, 1)}), confirmed against an isolated single-sprite diagnostic
tool kept permanently in the tree (\texttt{modules/renderers/sdl-gpu/examples/\allowbreak{}sdlgpu\_diag\_single\_sprite.cpp}, the same
``keep the diagnostic, don't throw it away once the bug is fixed'' precedent
Chapter~\ref{ch:direct3d-backends} already covers for \texttt{cna\_diag\_d3d12\_swapchain}).

\section{SkinnedEffect's storage-buffer workaround: a real, undocumented push-uniform cap}

Wiring \cnaclass{SkinnedEffect} through this renderer surfaced a genuine hardware/driver
constraint rather than a CNA design choice: \texttt{SDL\_PushGPUVertexUniformData}, the
mechanism every other stock effect above uses to deliver its per-draw uniforms, has a real,
undocumented cap of roughly 4096 bytes per slot on this renderer's own Vulkan driver --- found by
a real empirical binary-search spike, not read out of any SDL3 header or release note, since
the limit is not documented anywhere in \texttt{SDL\_gpu.h} itself. A 72-bone skin palette
(\(72 \times 64\) bytes \(= 4608\) bytes) exceeds that cap outright, so
\cnaclass{SkinnedEffect}'s bone matrices are uploaded through a real GPU storage buffer
(\texttt{SDL\_BindGPUVertexStorageBuffers}) instead of a uniform push --- every other stock
effect's smaller per-draw payloads (world/view/projection, \texttt{DiffuseColor}, alpha-test
parameters, light parameters) stay on the ordinary uniform-push path, since none of them come
close to the cap. This is exactly the kind of real, hardware-shaped constraint this book asks
every renderer chapter to name precisely rather than paper over with ``uniforms are uploaded
per draw'' as if the mechanism were uniform across every payload size.

\section{Custom ShaderEffect: a real runtime compile, built without the dev package installed}

Custom \cnaclass{ShaderEffect} support (Chapter~\ref{ch:shader-fx-gap}) is real on this renderer
too, and its own implementation is worth knowing about for a reason beyond the feature itself.
Since \texttt{SDL\_gpu} accepts only precompiled bytecode, \cnaclass{SdlGpuEffectRenderer} bridges
\cnaclass{ShaderEffect}'s public GLSL-source contract to that requirement with a genuine runtime
\texttt{libshaderc} GLSL-to-SPIR-V compile, linked directly into the renderer binary rather than
merely invoked at build time. Reading \texttt{SdlGpuRenderer.cpp} directly shows a small,
telling detail: this dev environment has no \texttt{libshaderc-dev} package installed, so there
is no real \texttt{shaderc.h} to include --- the renderer's own source hand-declares the exact
handful of \texttt{shaderc\_*} C entry points it needs. These include
\texttt{shaderc\_\allowbreak{}compiler\_\allowbreak{}initialize},
\texttt{shaderc\_\allowbreak{}compile\_\allowbreak{}into\_\allowbreak{}spv}, and
\texttt{shaderc\_\allowbreak{}result\_\allowbreak{}get\_\allowbreak{}bytes}. Their
\texttt{extern "C"} function signatures, shader-kind values, and optimization-level values are
hand-authored as named constants rather than pulled from a missing header.
This is a real, working engineering accommodation to a genuinely constrained dev environment,
not a shortcut around the feature itself --- the compiled shader still goes through the real
linked \texttt{libshaderc.so}, only the header describing its C ABI is hand-authored instead of
vendored. Building this custom-shader path also surfaced a genuine finding one layer over: the
existing \cnaclass{VulkanRenderer}'s own \cnaclass{ShaderEffect} support does not
actually runtime-compile GLSL text at all, despite \cnaclass{ShaderEffect}'s own documented
contract implying every renderer does --- a real, previously-unflagged discrepancy between one
renderer's implementation and the class's own stated cross-renderer promise, left for
Chapter~\ref{ch:vulkan-backend} to account for on its own terms rather than resolved here.

\subsection{The \texttt{.cnj} fixture exposes the remaining shader-compatibility boundary}

The runtime compiler does not make every GLSL text accepted elsewhere in CNA portable to this
renderer. The real \texttt{CnjEffectTest.LoadsRealCnjFixture} content test writes an ordinary
\texttt{.cnj} effect manifest plus a GLSL ES~3.00 vertex/fragment pair: \texttt{\#version 300 es},
unqualified stage outputs/inputs, and loose \texttt{uniform mat4 projection} /
\texttt{uniform sampler2D texture1} declarations. Under \texttt{SDL\_GPU}, that fixture fails
before a pixel is drawn. Its source is not in the SPIR-V-compatible dialect required by the
runtime compiler and by the resource counts CNA supplies to \texttt{SDL\_CreateGPUShader}.

The passing \texttt{SdlGpu\_ShaderEffect} regression makes the required contrast concrete. Its
two sources use desktop \texttt{\#version 450}, explicit \texttt{layout(location=...)} qualifiers
for every vertex attribute and stage I/O, and the exact fixed binding spaces this renderer
declares: the fragment sampler is \texttt{layout(set=2, binding=0)}, while the per-draw values
live in a named uniform block at \texttt{set=1} for the vertex stage and \texttt{set=3} for the
fragment stage. It compiles both stages at runtime, draws a white texture into a
\cnaclass{RenderTarget2D}, and reads back the effect's own blue-tinted value; repeating the
same draw without the effect must read white. Thus the supported dialect and binding contract
are proven by pixels, not merely by \texttt{IsEffectValid()}.

This is deliberately an open compatibility gap, not evidence that \cnaclass{ShaderEffect} is a
no-op on SDL\_GPU. Closing it needs a separately designed decision: either migrate the shared
\texttt{.cnj} fixture/source convention to the stricter portable dialect, or teach this renderer
to translate its existing loose GLSL into the blocks, locations, and bindings SDL\_gpu requires.
Neither is a safe one-line relaxation; the first changes content expectations across renderers,
and the second is a real shader-translation layer. The chapter therefore records the narrower,
honest current contract: custom effects work when authored for SDL\_GPU's binding dialect, while
the generic \texttt{.cnj} fixture is not yet cross-renderer portable to it.

\section{Pipeline state is snapshotted at draw time, except where it is deliberately not}

SDL\_gpu graphics pipelines are immutable enough that CNA cannot simply bind an XNA state object
at the end of a deferred frame and hope it represents every earlier draw. Each queued sprite or
3D draw therefore captures a \texttt{RenderStateSnapshot}: blend factors/functions, cull and
fill mode, stencil operations/masks/reference, and the ordinary depth-test fields travel with
the command that was issued under them. Pipeline lookup folds the pipeline-relevant portion into
a \texttt{size\_t} hash together with primitive topology, target color format, and sample count.
This replaced an earlier hand-packed integer key, whose fixed bit budget would have become
silently collision-prone as state dimensions were added.

The detail that keeps the cache both correct and finite is conditional hashing. Disabled blend
does not hash its six blend factors/functions; disabled stencil does not hash operations or masks;
and disabled two-sided stencil does not hash the counter-clockwise fields. Two commands that
produce the same pipeline consequently share it even if their irrelevant XNA state-object fields
differ, while a meaningful state change cannot reuse a pipeline with stale SDL\_gpu descriptors.

\begin{sourcenote}
The target format included in that key is not the public backbuffer-format field. Each frame
queries SDL for the swapchain texture's actual format and builds the matching pipeline.
Independently, construction chooses D24S8 or D32S8 by a real device-support query and later allocates that
combined attachment whenever a backbuffer frame needs it. A public
\texttt{DepthFormat::None} does not suppress it; absence means the device supported neither
combined format. Fullscreen remains an SDL-window operation outside this renderer hook, and
swapchain acquisition follows the resulting window without selecting formats from the public
enums.
\end{sourcenote}
Stencil reference itself is intentionally \emph{not} pipeline state: the captured value is sent
per draw through \texttt{SDL\_SetGPUStencilReference()}, the missing call that an adversarial
pixel test once exposed. A left-half stencil write followed by an Equal-reference reveal now
proves that this dynamic value reaches the GPU rather than merely being stored in CNA.

Viewport, scissor, and blend constants follow the same per-command rule. Every queued draw
reference snapshots the active viewport rectangle/depth range, scissor rectangle and enable bit,
and blend factor; replay applies them through \texttt{SDL\_SetGPUViewport},
\texttt{SDL\_SetGPUScissor}, and \texttt{SDL\_SetGPUBlendConstants} immediately before that
draw. Enabled scissors are clamped to the target; disabled or degenerate ones expand to its full
extent. The shared deferred-viewport/scissor suites queue several changes without a flush and
cover backbuffer, 2D-target, and cube-face routes, so a final live state cannot satisfy them by
accident.

Depth bias is pipeline-static in SDL\_gpu 3.5 rather than dynamic. CNA normalizes both
\cnaclass{RasterizerState.DepthBias} and \texttt{SlopeScaleDepthBias} by primitive topology,
includes their bits in every affected pipeline key, and fills SDL's constant and slope factors
when creating the pipeline. The dedicated test combines structural key checks, differential
depth pixels, and validation-fatal execution. Samplers likewise carry the full filter ordinal,
U/V address modes, and clamped \texttt{MaxAnisotropy} through both cache key and native
descriptor; anisotropy is no longer merely a stored capability bit.

\section{Deferred rendering still preserves the game's draw order}

Deferring all SDL\_gpu work until \texttt{Present()} once introduced a correctness bug hidden by
the renderer's family-specific command vectors: it rendered every 3D family first and every sprite
last, regardless of the order the game issued them. A background sprite, 3D model, and HUD sprite
therefore became background/model/HUD only by accident; even the opposite sprite--then--3D order
was replayed as 3D--sprite. This is visible whenever later opaque or blended content should cover
earlier content, not merely a performance detail.

Every \texttt{Queue*Draw()} and \texttt{QueueSprite()} call now appends a compact kind/index
reference beside its family-specific command. At render time,
\texttt{RenderQueuedDraws()} makes one chronological pass over those references and dispatches
the existing per-family issue routine only when its command belongs to the current target. No
sort is needed: append order is already the public draw-call order. Pipeline rebind avoidance is
then tracked across all kinds rather than only sprites, so fixing correctness does not require
forgetting the ordinary bind cache.

\texttt{SdlGpu\_DrawOrder} proves both directions with no alpha arithmetic ambiguity. It renders
an opaque red full-target sprite and an opaque green full-target 3D quad into separate targets
with depth disabled. Sprite then 3D must read back green; 3D then sprite must read back red.
The former fixed-family implementation produced red in both cases because sprites always ran
last, so the pair is a discriminating regression rather than a no-crash smoke check.

\section{Present timing belongs to the device manager, not a test-only renderer poke}

The SDL\_GPU tests establish a small but important rule about CNA's initialization lifecycle.
\cnaclass{GraphicsDeviceManager} defaults \texttt{SynchronizeWithVerticalRetrace} to true,
which is the right XNA-compatible default for an ordinary game. On this project's virtual test
display, however, there is no useful vertical-retrace signal; leaving it enabled made each frame
wait roughly one second. A 60- or 120-frame GPU regression test therefore consumed its CTest
budget without doing more useful rendering. The correct setup, used by all 22 current
SDL\_GPU examples and diagnostics, is public configuration made before
\texttt{Game::DoInitialize()} creates or resets the device:

\begin{lstlisting}
gdm_ = std::make_unique<GraphicsDeviceManager>(this);
gdm_->setPreferredBackBufferWidthProperty(64);
gdm_->setPreferredBackBufferHeightProperty(64);
gdm_->setSynchronizeWithVerticalRetraceProperty(false);
\end{lstlisting}

This order matters. Earlier tests called \texttt{SdlGpuRenderer::SetSwapInterval(0)}
directly as a local workaround. During initialization, though, the device manager converts its
still-true public property to \texttt{PresentInterval::One}; \cnaclass{GraphicsDevice::Reset}
then forwards that presentation interval to every renderer. The correct global forwarding was
itself a real repair: before it landed, the reset path did not call
\texttt{IGraphicsRenderer::SetSwapInterval()} at all. Once repaired, it rightly overwrote the
private pre-initialization poke, revealing the test timeout instead of preserving an accidental
ordering dependency.

For this renderer, interval zero asks SDL\_gpu for \texttt{SDL\_GPU\_PRESENTMODE\_IMMEDIATE},
falling back to \texttt{MAILBOX} if Immediate is unavailable; any positive interval selects
VSync. SDL\_gpu has no half-rate present mode, so XNA's \texttt{PresentInterval::Two} has the
same VSync outcome as \texttt{One}. That makes the example's explicit false a test-environment
choice, not a claim that games should globally disable synchronization: it uses the same public
property a game would use, at the phase of the lifecycle where CNA will not subsequently replace
it.

\section{Validation mode turned two tolerated violations into real fixes}

This renderer now passes SDL\_gpu's \texttt{debug\_mode} flag from CNA's ordinary build mode:
debug builds request validation and release builds do not. That deliberately modest toggle paid
for itself immediately. With validation previously hardcoded off, two API-contract violations
had appeared to work; with it on, each could hang the Vulkan driver rather than merely print a
warning. The lesson is more useful than a generic ``enable validation'' recommendation: a
renderer's release-mode pixels can look plausible while resource descriptions are still invalid.

First, the MSAA implementation for \cnaclass{RenderTargetCube} used a six-layer 2D-array
texture. SDL\_gpu forbids multisampling on an array texture. CNA now uses one single-layer,
multisampled 2D attachment for whichever cube face is currently rendering and resolves it into
the actual cube texture; cycling is permitted only for that temporary attachment, never for the
cube texture that must retain all six faces. Second, the initial automatic-mipmap implementation
for plain \cnaclass{Texture3D} and \cnaclass{TextureCube} requested only \texttt{SAMPLER} usage,
although SDL\_gpu's generator requires \texttt{COLOR\_TARGET} too. Widening the usage exposed a
third, deeper failure: SDL's Vulkan path created invalid 2D views and blits for depth planes that
no longer exist in smaller 3D mip levels.

That third finding is no longer open. \texttt{REMED-GFX-099} removed the invented automatic
generation behavior from plain \cnaclass{Texture3D}: XNA/FNA's contract allocates an authored
chain and lets each \texttt{SetData(level,...)} call populate the named level. The resource is
therefore \texttt{SAMPLER}-only again, with no target views or generator blits to violate the
Vulkan contract. A validation-fatal regression now covers exact level-zero and nonzero-mip
round trips, partial volumes, distinguishable depth planes, forward/reverse write order, NPOT
dimensions, and two alternating resources. Plain cube textures retain their separately tested
per-face behavior. This sequence is why validation must be treated as evidence with scope: it
first exposed tolerated violations, then forced the design back to the actual public resource
contract rather than merely suppressing their diagnostics.

\section{Render targets: a real use-after-free, closed before it ever shipped}

\cnaclass{RenderTarget2D} and \cnaclass{RenderTargetCube} support (with real multisampling and
mip regeneration) surfaced a genuine lifetime hazard while its own test was still being written,
not after: this renderer defers its actual render pass to \texttt{Present()} time, so a render
target destroyed \emph{before} that deferred pass executes was a real use-after-free the moment
anything still referenced its underlying GPU state. The fix,
\cnaclass{SdlGpuRenderTarget2DState} / \cnaclass{SdlGpuRenderTargetCubeState}, is a pair of
\texttt{shared\_ptr}-owned structs holding the actual GPU state independently of the
\cnaclass{SdlGpuRenderTargetRenderer} / \cnaclass{SdlGpuRenderTargetCubeRenderer} wrapper's own
C\texttt{++} object lifetime --- a short-lived, local render target destroyed mid-\texttt{Draw()}
now keeps its pending content alive long enough to render correctly and remain safely sampleable
elsewhere the same frame. The bug was confirmed via a real segfault in
\texttt{sdlgpu\_mrt\_test.cpp} before the fix landed, not merely reasoned about in the abstract.

\subsection{A regression that makes destruction precede submission}

The permanent regression test is deliberately stricter than ``a short-lived target does not
crash.'' In \texttt{modules/renderers/sdl-gpu/examples/\allowbreak{}sdlgpu\_rendertarget\_lifetime\_test.cpp}, one \(8\times8\)
\cnaclass{RenderTarget2D} is created as a local variable inside a single \texttt{Draw()} call,
cleared red, then sampled by a queued \cnaclass{SpriteBatch} draw into a still-live
\(16\times16\) destination target that was first cleared blue. The local variable reaches the
end of its scope before \texttt{Present()} invokes \texttt{EnsureFrameRendered()}, so the
sampled draw is still only a command waiting to be submitted:

\begin{lstlisting}
void DrawThroughShortLivedRenderTarget(GraphicsDevice& dev) {
    RenderTarget2D localRt(dev, 8, 8, false, SurfaceFormat::Color,
                           DepthFormat::None, 0, RenderTargetUsage::DiscardContents);
    dev.SetRenderTarget(&localRt);
    dev.Clear(Color::Red);
    dev.SetRenderTarget(nullptr);

    dev.SetRenderTarget(rtDest_.get());
    dev.Clear(Color::Blue);
    sb_->Begin();
    sb_->Draw(localRt, Rectangle(0, 0, 16, 16),
              Rectangle(0, 0, 8, 8), Color::White);
    sb_->End();
    dev.SetRenderTarget(nullptr);
} // localRt is destroyed here, before this frame is submitted
\end{lstlisting}

On the first frame the test reads the centre pixel of \texttt{rtDest\_} back and requires exact
red (within its small channel tolerance), proving that the queued sampling operation really
ran; a no-crash-only assertion could otherwise pass after silently omitting the draw. It then
performs the same create--clear--sample--destroy pattern for 120 frames, which checks that the
deferred-release queue is drained after each submission rather than merely retaining every
texture until shutdown. The implementation has two necessary lifetime stages: queued targets
and draw targets retain a \texttt{shared\_ptr} to their internal state until frame rendering has
finished, then that state's destructor places its raw color, depth, and MSAA handles in
\texttt{pendingTextureReleases\_}. After a successful command-buffer submission, the renderer
releases each queued handle through \texttt{SDL\_ReleaseGPUTexture()}. SDL may still fence the
physical memory internally; CNA's narrower rule is that a recording cannot release a handle
which an unsubmitted command still names.

\section{MRT is a real shader capability, not merely a multi-clear convenience}

\cnaclass{GraphicsDevice::SetRenderTargets()} makes the first
\cnaclass{RenderTarget2D} the primary target and records every later target as one simultaneous
SDL\_gpu color attachment. At \texttt{Present()} the primary's render pass contains
\(1+\lvert\texttt{mrtSiblings}\rvert\) \texttt{SDL\_GPUColorTargetInfo} entries; the custom
effect pipeline cache includes that attachment count, so a pipeline created for one color output
cannot accidentally be rebound for two. All attachments receive the same CNA blend state, just
as one \cnaclass{GraphicsDevice::BlendState} governs the whole draw rather than one state per
target.

That does \emph{not} mean every draw writes every attachment. CNA's stock sprite and 3D fragment
shaders declare one output, so their real, deliberately narrow MRT contract is: target zero
receives the draw while targets one onward can still be bound and cleared independently. The
second half of \texttt{sdlgpu\_mrt\_test.cpp} proves the broader custom-effect path by binding
two fresh targets and issuing exactly one \cnaclass{SpriteBatch} draw with a \texttt{\#version
450} fragment shader containing \texttt{layout(location=0) out vec4 outColorA} and
\texttt{layout(location=1) out vec4 outColorB}. Sampling a white texture with tint
\((0.2,0.4,0.8,1)\) must read \((51,102,204,255)\) from A and its shader-computed channel swap
\((102,204,51,255)\) from B. Those distinct readbacks rule out both an extra clear and a copied
single-output result: one fragment invocation really wrote two simultaneous attachments.

MRT lifetime follows the same ownership repair as ordinary deferred targets. Each immutable
\texttt{PassSegment} retains its primary and every secondary attachment as
\texttt{shared\_ptr<\allowbreak{}SdlGpuRenderTarget2DState>}; explicit clears operate on the segment rather
than a parallel wrapper-pointer list. Destroying a public secondary wrapper therefore cannot
invalidate an already-recorded bind cycle, and the render-target-lifetime regression covers the
create--bind--draw/clear--destroy--present sequence.

\section{Usage selects per-resource first-use load/store behavior}

Each 2D target and each cube face begins with independent first-use colour/depth/stencil flags.
The shared Discard bind records the FNA-style black/max-depth/zero-stencil clear into its new
bind-cycle segment; Preserve records none. Once a pass has written the resource, later cycles
load its stored aspects unless an explicit ordered clear says otherwise. MRT applies the
primary cycle's explicit colour clear to all simultaneous colour attachments while retaining
each attachment's owning state. The depth/stencil usage fixture renders distinct occlusion and
stencil-gate outcomes across unbind/rebind cycles, so the result is not inferred from colour.

Cube usage now reaches construction too. A single-sample face naturally loads its cube layer;
a multisampled face owns a persistent per-face attachment and uses
\texttt{RESOLVE\_AND\_STORE} for Preserve so the next cycle can load its samples. Discard may
use plain \texttt{RESOLVE} when no later segment in the frame needs them. The six-face MSAA/mip
and depth/stencil usage suites cover these paths. \texttt{PlatformContents} follows the shared
FNA predicate \texttt{usage != DiscardContents} and therefore preserves like Preserve; cube
and depth/stencil usage fixtures assert that policy explicitly. See
\S\ref{sec:backend-rendertarget-usage-matrix}.

\subsection{A worked example: \cnaclass{Texture3D}'s own \texttt{cycle=true} orphaned-write bug}

\cnaclass{Texture3D} support on this renderer --- sub-volume upload/readback plus mipmap generation
--- surfaced a real data-corruption bug of exactly the kind this book keeps returning to: code that compiles cleanly, throws
nothing, and produces a wrong answer silently. The relevant upload operation is
\texttt{SDL\_UploadToGPUTexture}. Its \texttt{cycle} parameter offers a real performance
trade-off SDL3's own documentation states plainly: \texttt{cycle=true} swaps the texture to a
fresh, separate underlying GPU resource on each call, avoiding a stall if the GPU is still
reading the previous resource's content. That is correct and desirable for \cnaclass{Texture2D}'s
single-full-texture-replace \texttt{SetData()} path. CNA's SDL\_GPU texture renderer calls that
path \texttt{UpdatePixels}; it deliberately retains \texttt{cycle=true}. The first
\cnaclass{Texture3D} implementation copied that choice without accounting for a 3D texture's
multiple writes.
The assumption breaks for \cnaclass{Texture3D} specifically, because its content is built up
through \emph{multiple independent} sub-volume and per-level \texttt{SetData()} calls that all
need to land on the \emph{same} underlying resource --- with \texttt{cycle=true}, an earlier
partial write (say, a sub-volume at $z=0$) silently orphaned itself onto an abandoned GPU
resource the moment a second, later write (a sub-volume at $z=1$) followed it, reading back as
zero/uninitialized rather than the real content actually uploaded.

The real regression test that caught this (\texttt{modules/renderers/sdl-gpu/examples/\allowbreak{}sdlgpu\_texture3d\_test.cpp})
is deliberately shaped to make the bug unmissable rather than merely possible to catch:

\begin{lstlisting}
Texture3D tex(dev, 4, 4, 2, false, SurfaceFormat::Color);
// A 2x2x2 sub-volume at offset (1,1,0) -- red at z=0, green at z=1 -- deliberately
// off-center within the full 4x4x2 volume, and two SEPARATE SetData() calls, one per Z slice.
tex.SetData(0, /*x=*/1, /*y=*/1, /*z=*/0, 2, 2, 1, redPixels.data(), redPixels.size());
tex.SetData(0, /*x=*/1, /*y=*/1, /*z=*/1, 2, 2, 1, greenPixels.data(), greenPixels.size());
// With cycle=true, the z=0 write above would already be silently orphaned by the time this
// GetData() runs, reading back zero instead of red.
std::vector<Color> gotRed(4), gotGreen(4);
tex.GetData(0, 1, 1, 0, 2, 2, 1, gotRed.data(), gotRed.size() * sizeof(Color));
tex.GetData(0, 1, 1, 1, 2, 2, 1, gotGreen.data(), gotGreen.size() * sizeof(Color));
\end{lstlisting}

Two design choices make this test genuinely discriminating rather than accidentally passing
either way: the sub-volume is off-center (a bug that only corrupted, say, the texture's
origin corner would still show up), and each Z slice gets its own distinct solid color rather
than reusing one color for the whole volume --- a uniform-color test cannot tell correctly
-downsampled or correctly-retained content apart from an out-of-range read that happens to
land on more of that same color, the identical discriminating-test principle
Section~\ref{sec:sdlgpu-what-remains-open}'s own generated-mipmap gap needs applied to close it.
The fix, confirmed against the real
\texttt{SdlGpuTexture3DRenderer::\allowbreak{}SetData} implementation
directly, is a one-word change: pass \texttt{cycle=false} for every \cnaclass{Texture3D}
upload, leaving \cnaclass{Texture2D}'s own \texttt{cycle=true} single-replace path untouched
and correct as-is. The identical fix (and the identical reasoning, cited by name in that
renderer class's own source comment) was applied proactively to
\cnaclass{TextureCube}'s per-face uploads too, since a cube map is built up through six
separate per-face \texttt{SetData()} calls that must all land on the same underlying resource --- caught
and fixed before shipping rather than found by reproducing the same bug a second time.

\section{Forced swapchain failure preserves the queued frame}

SDL\_gpu distinguishes two acquisition outcomes that a renderer must not conflate. A successful
\texttt{SDL\_WaitAndAcquireGPUSwapchainTexture()} call with a null texture is a documented,
non-error case such as a minimized window: CNA submits the command buffer and simply skips that
frame. A false return is a hard acquisition failure. SDL forbids cancelling the command buffer
after this call, so CNA submits it first and then throws with the captured SDL error. Crucially,
it leaves \texttt{framePending\_} true: a caller which recovers the window must be able to submit
the same queued clear and draws rather than silently lose them.

\texttt{sdlgpu\_\allowbreak swapchain\_\allowbreak recovery\_\allowbreak test.cpp} proves that latter path without pretending a real
device loss is easy to reproduce. On frame 10 it calls SDL's own
\texttt{SDL\_\allowbreak ReleaseWindow\allowbreak FromGPU\allowbreak Device()}, makes \texttt{Present()} fail, then
uses SDL's matching claim call to reclaim the identical window before calling \texttt{Present()}
again. The test
requires all five stages: nine ordinary frames, the induced throw, successful reclaim, successful
presentation of the preserved frame, and 30 further frames without exception. This is a bounded,
controlled proof of recovery after a hard acquisition failure; it does not claim to solve the
separate lazy-proxy backbuffer readback path described below.

\section{What remains open}
\label{sec:sdlgpu-what-remains-open}

Three limits are worth stating precisely rather than glossing over. First, occlusion queries are a
\emph{permanent} limitation of the underlying API, not a task gap: the vendored
\texttt{SDL\_gpu.h} has no query-pool or occlusion-query type anywhere in it (only
\texttt{SDL\_GPUFence} for CPU/GPU synchronization), so \texttt{CreateOcclusionQuery()}
correctly returns \texttt{nullptr} on this renderer, the identical posture \texttt{HEADLESS} and
\texttt{SOFTWARE} already take for their own, unrelated reasons --- worth re-opening only if a
future SDL3 release adds real query support, not before. Second, the generic \texttt{.cnj} custom-effect fixture remains a
cross-renderer portability gap: its loose GLSL ES 3.00 declarations are not accepted by this
renderer's explicit-location, uniform-block binding contract, even though the dedicated
\texttt{\#version 450} effect regression is pixel-verified. That choice is deliberately left
open pending either a shared content-dialect migration or a real shader-translation layer.
Third, hardware instancing is absent at a different layer: this otherwise 3D-capable renderer
does not override \texttt{DrawInstancedPrimitivesEx()}, so every valid public
\texttt{DrawInstancedPrimitives()} call reaches the common
\texttt{std::runtime\_error} default. SDL\_gpu itself exposes per-instance vertex input; the
gap is CNA-side pipeline/binding work, not a permanent underlying-API limitation. The public
\cnaclass{GraphicsCapability::Instancing} entry now exists, but this renderer does not override
the permissive base policy and therefore incorrectly reports it as supported. This is a
capability overclaim: callers cannot safely use the bit as a guard until SDL\_GPU either
implements the draw route or returns false explicitly
(Section~\ref{sec:drawex-instancing-matrix}).

Five former entries no longer belong in this list. Depth bias and slope-scale bias are
normalized per polygon topology, hashed into every relevant pipeline key, and applied through
SDL's pipeline-static rasterizer fields; a dedicated pixel/validation test covers them.
\texttt{MaxAnisotropy} is part of the complete sampler key and native descriptor. Secondary MRT
attachments are retained as \texttt{shared\_ptr}s in the immutable pass segment, closing the raw
pointer lifetime hazard. Finally, ordinary \cnaclass{Texture2D} allocates its declared native
mip count and overrides \texttt{UpdatePixelsLevel}; the native-allocation/upload regression
proves that a nonzero level reaches the GPU resource. None of these closures adds automatic
downsampling to plain \cnaclass{Texture2D}: its upper levels remain explicitly authored, as in
the XNA/FNA contract.

Backbuffer readback is also closed without charging games that never use it. SDL's swapchain
texture remains permanently write-only, but the first \texttt{GetBackBufferData()} lazily enables
a self-owned, swapchain-sized proxy. The pending frame renders into that proxy, one 1:1 blit
presents it, and \texttt{ReadBackbuffer()} downloads the requested logical rectangle through a
transfer buffer and fence, converting BGRA to public RGBA when needed. Before the first read the
proxy and blit do not exist. First-read, whole-logical-dimension, and fully-written-range tests
cover the observable contract.

Beyond these three remaining limits, Windows (D3D12 driver) and macOS/iOS (Metal driver)
remain code paths only, not validation claims, until run on real (or Wine/DXVK-class) hardware
--- the same phased-claim discipline this book already applies to \texttt{D3D11} /
\texttt{D3D12} (Chapter~\ref{ch:direct3d-backends}) and \texttt{WebGPU}
(Chapter~\ref{ch:webgpu-backend}).

\section{Clear state is per segment and issue-ordered}

This renderer carries independent colour, depth, and stencil requests on each backbuffer,
2D-target/MRT, or cube-face segment. Consecutive clears before any draw coalesce safely into the
segment's load action, aspect by aspect. Once a draw has made ordering observable, the next
clear opens another segment over the same destination; the following native pass loads
unselected aspects, clears selected ones, and then replays later draws. MSAA segments use
\texttt{RESOLVE\_AND\_STORE} whenever a later segment must load their samples.

The renderer-neutral clear-options and ordered-clear fixtures distinguish isolated aspects,
draw--clear versus clear--draw, multiple clears, and backbuffer/2D/cube/MRT destinations without
an intermediate flush. Pass-boundary and backbuffer-order suites separately prove that target
rebinding and target--backbuffer interleaving remain in public order. Deferred submission is
therefore an implementation strategy, not a remaining clear-order exception; the comparative
matrix is \S\ref{sec:backend-clear-matrix}.

\section{Verification methodology}

The executable claims above have dedicated tests rather than resting on code reads alone.
Counting the actual module-local registrations directly
(\texttt{modules/renderers/sdl-gpu/\allowbreak{}examples/CMakeLists.txt}) rather than trusting an old plan
snapshot shows \textbf{83} \texttt{SdlGpu\_*} CTest registrations today, spanning the 2D and 3D
smoke tests, every stock effect (including \cnaclass{EnvironmentMapEffect},
\cnaclass{SkinnedEffect}, and the \texttt{CNAEXT} \cnaclass{PbrEffect} this book's WebGPU chapter
already introduced), custom \cnaclass{ShaderEffect}, every render-target variant (2D, cube,
multiple render targets, MSAA), both extra texture types, sampler/render state, draw order, and
swapchain recovery, plus the render-target-lifetime regression above --- run, like every other renderer's own CTest suite, against this project's
own \texttt{cna\_register\_renderer\_test()} convention, gracefully skipping with a clear reason
rather than hard-failing when no real \texttt{DISPLAY} / \texttt{WAYLAND\_DISPLAY} is present, the
identical convention \texttt{D3D11} / \texttt{D3D12} / \texttt{WebGPU} already follow. All verification
to date is against this project's own real dev-machine GPU via \texttt{SDL\_gpu}'s Vulkan
driver, the only driver this project's own prebuilt SDL3 packages compile in on Linux --- the
same single-driver-verified, other-drivers-unverified posture this chapter's own ``What remains
open'' section already names honestly rather than implying broader coverage than actually
exists.

